\ifdefined\gkver\else\def\gkver{student}\fi
\documentclass[\gkver]{gaokaozhenti}
\begin{document}

\chapter{2007年广东理科基础}
%% paperType: 理科基础物理部分

\begin{enumerate}
\item
%% number: 1
%% paperName: 2007年广东理科基础第 1 题 2 分
%% typeId: 1
%% type: 单选题
%% chapter: 功和能
%% point: 功
%% method: 无
%% score: 2
%% degree: 950
%% duplicateId: 0
%% body:
下列物理量为标量的是\xzanswer{D}

\fourchoices[answer=D]
{平均速度}
{加速度}
{位移}
{功}

%% answer: D
%% memo:
\memoanswer{本题原卷及材料中未提供详解，待补充。}

\item
%% number: 2
%% paperName: 2007年广东理科基础第 2 题 2 分
%% typeId: 1
%% type: 单选题
%% chapter: 直线运动
%% point: 自由落体运动
%% method: 无
%% score: 2
%% degree: 950
%% duplicateId: 0
%% body:
关于自由落体运动，下列说法正确的是\xzanswer{C}

\fourchoices[answer=C]
{物体竖直向下的运动就是自由落体运动}
{加速度等于重力加速度的运动就是自由落体运动}
{在自由落体运动过程中，不同质量的物体运动规律相同}
{物体做自由落体运动位移与时间成反比}

%% answer: C
%% memo:
\memoanswer{本题原卷及材料中未提供详解，待补充。}

\item
%% number: 3
%% paperName: 2007年广东理科基础第 3 题 2 分
%% typeId: 1
%% type: 单选题
%% chapter: 直线运动
%% point: 运动图象
%% method: 无
%% score: 2
%% degree: 850
%% duplicateId: 0
%% body:
如图所示是某物体做直线运动的速度图象，下列有关物体运动情况判断正确的是\xzanswer{A}

\onepicture[w=5cm]{figs/03.png}

\fourchoices[answer=A]
{前2秒加速度为$5\Umsq$}
{4 s 末物体回到出发点}
{6 s 末物体距出发点最远}
{8 s 末物体距出发点最远}

%% answer: A
%% memo:
\memoanswer{本题原卷及材料中未提供详解，待补充。}

\item
%% number: 4
%% paperName: 2007年广东理科基础第 4 题 2 分
%% typeId: 1
%% type: 单选题
%% chapter: 力物体的平衡
%% point: 多力平衡
%% method: 无
%% score: 2
%% degree: 850
%% duplicateId: 0
%% body:
受斜向上的恒定拉力作用，物体在粗糙水平面上做匀速直线运动，则下列说法正确的是\xzanswer{D}

\fourchoices[answer=D]
{拉力在竖直方向的分量一定大于重力}
{拉力在竖直方向的分量一定等于重力}
{拉力在水平方向的分量一定大于摩擦力}
{拉力在水平方向的分量一定等于摩擦力}

%% answer: D
%% memo:
\memoanswer{本题原卷及材料中未提供详解，待补充。}

\item
%% number: 5
%% paperName: 2007年广东理科基础第 5 题 2 分
%% typeId: 1
%% type: 单选题
%% chapter: 曲线运动
%% point: 曲线运动条件
%% method: 无
%% score: 2
%% degree: 850
%% duplicateId: 0
%% body:
质点在一平面内沿曲线由$P$运动到$Q$，如果$v$、$a$、$F$分别表示质点运动过程中的速度、加速度和受到的合外力，下列图象可能正确的是\xzanswer{D}

\fourchoices[answer=D, ispicture=true, w=6.5cm]
{figs/05a.png}
{figs/05b.png}
{figs/05c.png}
{figs/05d.png}

%% answer: D
%% memo:
\memoanswer{本题原卷及材料中未提供详解，待补充。}

\item
%% number: 6
%% paperName: 2007年广东理科基础第 6 题 2 分
%% typeId: 1
%% type: 单选题
%% chapter: 曲线运动
%% point: 平抛运动
%% method: 无
%% score: 2
%% degree: 850
%% duplicateId: 0
%% body:
质点从同一高度水平抛出，不计空气阻力，下列说法正确的是\xzanswer{D}

\fourchoices[answer=D]
{质量越大，水平位移越大}
{初速度越大，落地时竖直方向速度越大}
{初速度越大，空中运动时间越长}
{初速度越大，落地速度越大}

%% answer: D
%% memo:
\memoanswer{本题原卷及材料中未提供详解，待补充。}

\item
%% number: 7
%% paperName: 2007年广东理科基础第 7 题 2 分
%% typeId: 1
%% type: 单选题
%% chapter: 功和能
%% point: 功
%% method: 无
%% score: 2
%% degree: 750
%% duplicateId: 0
%% body:
人骑自行车下坡，坡长$l=500\Um$，坡高$h=8\Um$，人和车总质量为$100\Ukg$，下坡时初速度为$4\Ums$，人不踏车的情况下，到达坡底时车速为$10\Ums$，$g$取$10\Umsq$，则下坡过程中阻力所做的功为\xzanswer{B}

\fourchoices[answer=B]
{$-400\UJ$}
{$-3800\UJ$}
{$-50000\UJ$}
{$-4200\UJ$}

%% answer: B
%% memo:
\memoanswer{本题原卷及材料中未提供详解，待补充。}

\item
%% number: 8
%% paperName: 2007年广东理科基础第 8 题 2 分
%% typeId: 1
%% type: 单选题
%% chapter: 曲线运动
%% point: 竖直平面圆周运动
%% method: 无
%% score: 2
%% degree: 850
%% duplicateId: 0
%% body:
游客乘坐过山车，在圆弧轨道最低点处获得的向心加速度达到$20\Umsq$，$g$取$10\Umsq$，那么此位置座椅对游客的作用力相当于游客重力的\xzanswer{C}

\fourchoices[answer=C]
{1倍}
{2倍}
{3倍}
{4倍}

%% answer: C
%% memo:
\memoanswer{本题原卷及材料中未提供详解，待补充。}

\item
%% number: 9
%% paperName: 2007年广东理科基础第 9 题 2 分
%% typeId: 1
%% type: 单选题
%% chapter: 功和能
%% point: 功
%% method: 无
%% score: 2
%% degree: 850
%% duplicateId: 0
%% body:
一人乘电梯从1楼到30楼，在此过程中经历了先加速、后匀速、再减速的运动过程，则电梯支持力对人做功情况是\xzanswer{D}

\fourchoices[answer=D]
{加速时做正功，匀速时不做功，减速时做负功}
{加速时做正功，匀速和减速时做负功}
{加速和匀速时做正功，减速时做负功}
{始终做正功}

%% answer: D
%% memo:
\memoanswer{本题原卷及材料中未提供详解，待补充。}

\item
%% number: 10
%% paperName: 2007年广东理科基础第 10 题 2 分
%% typeId: 1
%% type: 单选题
%% chapter: 功和能
%% point: DIS研究机械能守恒定律
%% method: 无
%% score: 2
%% degree: 950
%% duplicateId: 0
%% body:
某位同学做“验证机械能守恒定律”的实验，下列操作步骤中错误的是\xzanswer{C}

\fourchoices[answer=C]
{把打点计时器固定在铁架台上，用导线连接到低压交流电源}
{将连有重锤的纸带过限位孔，将纸带和重锤提升到一定高度}
{先释放纸带，再接通电源}
{更换纸带，重复实验，根据记录处理数据}

%% answer: C
%% memo:
\memoanswer{本题原卷及材料中未提供详解，待补充。}

\item
%% number: 11
%% paperName: 2007年广东理科基础第 11 题 2 分
%% typeId: 1
%% type: 单选题
%% chapter: 曲线运动
%% point: 人造地球卫星
%% method: 无
%% score: 2
%% degree: 650
%% duplicateId: 0
%% body:
现有两颗绕地球匀速圆周运动的人造地球卫星 A 和 B，它们的轨道半径分别为$r_{A}$和$r_{B}$。如果$r_{A}>r_{B}$，则\xzanswer{A}

\fourchoices[answer=A]
{卫星 A 的运动周期比卫星 B 的运动周期大}
{卫星 A 的线速度比卫星 B 的线速度大}
{卫星 A 的角速度比卫星 B 的角速度大}
{卫星 A 的加速度比卫星 B 的加速度大}

%% answer: A
%% memo:
\memoanswer{本题原卷及材料中未提供详解，待补充。}

\item
%% number: 12
%% paperName: 2007年广东理科基础第 12 题 2 分
%% typeId: 1
%% type: 单选题
%% chapter: 电场
%% point: 电场线
%% method: 无
%% score: 2
%% degree: 850
%% duplicateId: 0
%% body:
如图所示，实线是一簇未标明方向的由点电荷$Q$产生的电场线，若带电粒子$q$（$|Q|\gg|q|$）由a运动到b，电场力做正功。已知在a、b两点粒子所受电场力分别为$F_{a}$、$F_{b}$，则下列判断正确的是\xzanswer{A}

\onepicture[w=4.5cm]{figs/12.png}

\fourchoices[answer=A]
{若$Q$为正电荷，则$q$带正电，$F_{a}>F_{b}$}
{若$Q$为正电荷，则$q$带正电，$F_{a}<F_{b}$}
{若$Q$为负电荷，则$q$带正电，$F_{a}>F_{b}$}
{若$Q$为负电荷，则$q$带正电，$F_{a}<F_{b}$}

%% answer: A
%% memo:
\memoanswer{本题原卷及材料中未提供详解，待补充。}

\item
%% number: 13
%% paperName: 2007年广东理科基础第 13 题 2 分
%% typeId: 1
%% type: 单选题
%% chapter: 电路
%% point: 电容
%% method: 无
%% score: 2
%% degree: 850
%% duplicateId: 0
%% body:
电容器是一种常用的电子元件。对电容器认识正确的是\xzanswer{A}

\fourchoices[answer=A]
{电容器的电容表示其储存电荷能力}
{电容器的电容与它所带的电量成正比}
{电容器的电容与它两极板间的电压成正比}
{电容的常用单位有$\UmuF$和$\UpF$，$1\UMF=10^{3}\UpF$}

%% answer: A
%% memo:
\memoanswer{本题原卷及材料中未提供详解，待补充。}

\item
%% number: 14
%% paperName: 2007年广东理科基础第 14 题 2 分
%% typeId: 1
%% type: 单选题
%% chapter: 电路
%% point: 故障分析
%% method: 无
%% score: 2
%% degree: 850
%% duplicateId: 0
%% body:
用电压表检查图示电路中的故障，测得$U_{ad}=5.0\UV$，$U_{ed}=0\UV$，$U_{be}=0\UV$，$U_{ab}=5.0\UV$，则此故障可能是\xzanswer{B}

\onepicture[w=4.5cm]{figs/14.png}

\fourchoices[answer=B]
{$L$ 断路}
{$R$ 断路}
{$R'$ 断路}
{S 断路}

%% answer: B
%% memo:
\memoanswer{本题原卷及材料中未提供详解，待补充。}

\item
%% number: 15
%% paperName: 2007年广东理科基础第 15 题 2 分
%% typeId: 1
%% type: 单选题
%% chapter: 电路
%% point: 伏安法测电阻
%% method: 无
%% score: 2
%% degree: 750
%% duplicateId: 0
%% body:
如图所示为伏安法测电阻的一种常用电路。以下分析正确的是\xzanswer{A}

\onepicture[w=4cm]{figs/15.png}

\fourchoices[answer=A]
{此接法的测量值大于真实值}
{此接法的测量值小于真实值}
{此接法要求待测电阻值小于电流表内阻}
{开始实验时滑动变阻器滑动头$P$应处在最左端}

%% answer: A
%% memo:
\memoanswer{本题原卷及材料中未提供详解，待补充。}

\item
%% number: 16
%% paperName: 2007年广东理科基础第 16 题 2 分
%% typeId: 1
%% type: 单选题
%% chapter: 磁场
%% point: 磁感线
%% method: 无
%% score: 2
%% degree: 950
%% duplicateId: 0
%% body:
磁体之间的相互作用是通过磁场发生的。对磁场认识正确的是\xzanswer{C}

\fourchoices[answer=C]
{磁感线有可能出现相交的情况}
{磁感线总是由 N 极出发指向 S 极}
{某点磁场的方向与放在该点小磁针静止时 N 极所指方向一致}
{若在某区域内通电导线不受磁场力的作用，则该区域的磁感应强度一定为零}

%% answer: C
%% memo:
\memoanswer{本题原卷及材料中未提供详解，待补充。}

\item
%% number: 17
%% paperName: 2007年广东理科基础第 17 题 2 分
%% typeId: 1
%% type: 单选题
%% chapter: 磁场
%% point: 左手定则
%% method: 无
%% score: 2
%% degree: 850
%% duplicateId: 0
%% body:
如图，用两根相同的细绳水平悬挂一段均匀载流直导线$MN$，电流$I$方向从$M$到$N$，绳子的拉力均为$F$。为使$F=0$，可能达到要求的方法是\xzanswer{C}

\onepicture[w=5cm]{figs/17.png}

\fourchoices[answer=C]
{加水平向右的磁场}
{加水平向左的磁场}
{加垂直纸面向里的磁场}
{加垂直纸面向外的磁场}

%% answer: C
%% memo:
\memoanswer{本题原卷及材料中未提供详解，待补充。}

\item
%% number: 18
%% paperName: 2007年广东理科基础第 18 题 2 分
%% typeId: 1
%% type: 单选题
%% chapter: 磁场
%% point: 洛伦兹力
%% method: 无
%% score: 2
%% degree: 850
%% duplicateId: 0
%% body:
如图，在阴极射管正下方平行放置一根通有足够强直流电流的长直导线，且导线中电流方向水平向右，则阴极射线将会\xzanswer{A}

\onepicture[w=5cm]{figs/18.png}

\fourchoices[answer=A]
{向上偏转}
{向下偏转}
{向纸内偏转}
{向纸外偏转}

%% answer: A
%% memo:
\memoanswer{本题原卷及材料中未提供详解，待补充。}

\item
%% number: 19
%% paperName: 2007年广东理科基础第 19 题 2 分
%% typeId: 1
%% type: 单选题
%% chapter: 气体性质
%% point: 能源的开发与利用
%% method: 无
%% score: 2
%% degree: 950
%% duplicateId: 0
%% body:
下述做法能改善空气质量的是\xzanswer{B}

\fourchoices[answer=B]
{以煤等燃料作为主要生活燃料}
{利用太阳能、风能和氢能等能源替代化石能源}
{鼓励私人购买和使用汽车代替公交车}
{限制使用电动车}

%% answer: B
%% memo:
\memoanswer{本题原卷及材料中未提供详解，待补充。}

\end{enumerate}

\end{document}
